\cleardoublepage

\section{结论与展望}

\subsection{工作总结}

本文针对面向科学场景的智能体工具调用优化问题，从工具选择与任务拆解两个朴素 ReAct 范式中尚未被结构化处理的子问题出发，提出并实现了完整的算法、工程与评测体系。
具体而言，
\textbf{在算法层面}，提出了检索增强工具选择算法 RATS 与结构化 DAG 任务拆解算法 DAGPlanner，二者形式化清晰、可独立消融，并配合 4 条科学场景提示工程规则共同构成最终方案 \texttt{Ours\_full}；
\textbf{在工程层面}，构建了 12 工具的标准化科学计算工具库、6 种 Agent 策略的统一框架、覆盖三维判分的端到端评估框架；
\textbf{在实验层面}，在 9 配置 $\times$ 54 题共 486 次 run 的全量评测中验证了方法的有效性——\texttt{Ours} 与 \texttt{Ours\_full} 均达到 100\% 成功率，相比朴素 ReAct 提升 9.3 个百分点，相比业界最强基线 Reflexion 提升 5.6 个百分点。
更重要的是，本文方法在长依赖链 \texttt{composite} 类别上取得了独占的 100\% 成功率，并通过 DAGPlanner 把工具调用轨迹外显为可纯代码校验的 DAG 结构，在轨迹可解释性这一维度上相对自然语言 Plan-and-Solve 实现了质的提升。

\subsection{研究局限性}

囿于本科毕业设计的体量与时间，本文工作仍存在若干诚实需要承认的局限：

\textbf{第一，测试集规模偏小}。
54 题相对 ScienceAgentBench~\cite{chen2024scienceagentbench}（数百道任务）与 Berkeley Function Calling Leaderboard~\cite{bfcl2024}（千题级）等公开评测仍属中小规模；
本文虽然通过定向压力题部分缓解了这一问题，但外部可推广性仍需要更谨慎的表述，特别是 \texttt{Ours} 已在 54 题上饱和到 100\% 这一点限制了进一步消融的判别力。

\textbf{第二，评测仅限单基座模型}。
所有实验均基于通义千问 \texttt{qwen-plus-latest} 与 \texttt{text-embedding-v3} 完成，未在更强（如 GPT-4o）或更弱（如 7B 量级开源模型）的基座上交叉验证。
RATS 与 DAGPlanner 在不同基座下的相对优势是否保持，尚需后续实验支撑。

\textbf{第三，工具库规模偏小}。
12 个工具在功能多样性上虽然覆盖了数值计算、统计分析、可视化三大类常见科研需求，但对 RATS 的检索价值而言仍属“刚性显著但优势未能充分释放”的规模。
RATS 在 100+ 工具场景下的真正优势仍主要依赖外推，缺乏同口径的大规模工具库压力测试。

\textbf{第四，DAG 模式不支持循环与分支}。
DAGPlanner 的 Plan Schema 是有向无环图，天然适合链式、可拆分为明确子目标的科学计算流程；
但对需要显式循环（如“迭代直到收敛”）或强分支探索（如蒙特卡洛搜索）的策略类任务，DAG 表达力不足，需要切换到 ReAct 安全网或更通用的规划范式。

\textbf{第五，评测指标的检查器约束与智能体自主推理之间存在张力}。
评测使用三维判分可以高效量化结果，但也可能把“合理解法却未按检查器要求调用某工具”的情况判为失败，例如某些步骤在代数上可化简为开方而绕开 \texttt{equation\_solver}。
此类现象提醒我们：科研场景下检查器约束并非“科学真理性”，而是评估机制相对智能体自主推理的边界。

\subsection{未来工作}

围绕上述局限，未来可在以下五个方向继续推进：

\textbf{方向一：更大规模的工具库压力测试}。
将工具库扩展至 50--100 工具规模（覆盖更多统计、信号处理、生物信息学专项），重新跑 9 配置全量评测，验证 RATS 在大工具规模下的优势是否如预期反转。

\textbf{方向二：接入公开 Agent 评测基准}。
完成 ScienceAgentBench、BFCL v3、GAIA~\cite{mialon2023gaia} 等公开 benchmark 的接口对齐，把本文方法置于业界统一口径下做对比。

\textbf{方向三：跨基座模型评估}。
在 GPT-4o、Claude 3.5 Sonnet、DeepSeek-R1 等不同基座上重复主实验，量化 RATS 与 DAGPlanner 对基座的依赖度。

\textbf{方向四：DAGPlanner 表达力扩展}。
为 Plan Schema 增加循环步骤（\texttt{loop\_until} 与终止条件）与条件分支（\texttt{branch\_on}），并为校验器扩展对应的判定规则；
这一扩展将让 DAGPlanner 的适用边界从“链式确定性任务”拓宽到“迭代+分支的探索性任务”。

\textbf{方向五：RATS 重排权重的自动学习}。
当前权重 $(w_e, w_a, w_c, w_h)=(0.6, 0.2, 0.1, 0.1)$ 是经验值；
后续可基于历史成功率反馈，用强化学习或贝叶斯优化对权重做自动调优，进一步释放规则化重排的精度上限。

\subsection{结语}

“面向科学场景的智能体工具调用优化”是一个介于通用智能体研究与垂直领域工程之间的细分问题。
本文的贡献并不在于宣称一种全新的智能体范式，而在于把通用 ReAct 范式中的两个具体瓶颈——工具选择与任务拆解——在科学场景的精度约束下，分别用一个轻量、形式化、可独立消融的算法模块加以解决。
这条“在通用范式之上做轻量结构化”的路径，相比从零构建一个垂直领域的专用智能体，更易于复用、更便于消融、也更尊重现有大模型生态的成熟基础设施。
我们希望这一路径所体现的工程审慎与方法克制，能够为后来者在垂直领域智能体研究中提供一份可参考的实现范式。
